\chapter{Engineering Notes}
\label{app:eng-notes}

This appendix records implementation details that are needed to reproduce the
training runs but are not part of the method or the results. They are collected
here so that Chapter~\ref{ch:implementation} can stay focused on the components
that carry the contribution.

% --------------------------------------------------------------------------
\section{Library Workarounds}
\label{app:lib-workarounds}
% --------------------------------------------------------------------------

\textbf{TRL vLLM availability check.}
In TRL~0.14.0, \texttt{is\_vllm\_available()} returns the two-element tuple
\texttt{(False, None)}.
A non-empty tuple is truthy in Python, so the boolean guard inside
\texttt{GRPOTrainer} treats vLLM as available and attempts to initialise a vLLM
inference engine, which is not installed in this environment.
Overriding the helper to return \texttt{False} before the TRL import avoids the
spurious initialisation:
\begin{lstlisting}[caption={TRL vLLM check override}]
import trl.utils.import_utils as _iu
_iu.is_vllm_available = lambda: False
\end{lstlisting}
\texttt{GRPOConfig} additionally sets \texttt{use\_vllm=False}.

\textbf{Transformers KV-cache regression.}
With \texttt{gradient\_checkpointing} enabled and \texttt{model.train()} active on
a Qwen2 model under Transformers~5.x, \texttt{Qwen2DecoderLayer} forces
\texttt{past\_key\_values=None} on every forward pass, including during
autoregressive generation.
This disables the KV cache at generation time, so each completion token is
recomputed over the full sequence length.
Completions remain correct but generation slows substantially, which matters for a
loop that samples $G$ completions per prompt at every step.
A wrapper around \texttt{model.generate} that restores the correct generation-mode
context for the duration of the call avoids the regression.

% --------------------------------------------------------------------------
\section{Unattended-Run Tooling}
\label{app:run-tooling}
% --------------------------------------------------------------------------

\textbf{Health monitoring.}
A TRL callback fires once an hour and writes a health summary: an SSH ping to the
KCOV VM to confirm the reward oracle is reachable, the verdict distribution of the
last completed batch (ACCEPTED, REJECTED, encode-fail), and the current cumulative
PC count.
When training runs against a remote reward server the callback disables itself,
since the cloud session logs serve the same purpose.

\textbf{Checkpoint resume.}
A \texttt{--resume} flag scans the checkpoint directory, selects the latest
checkpoint, and passes it to \texttt{GRPOTrainer} as
\texttt{resume\_from\_checkpoint}, so restarting after an interruption requires no
manual checkpoint bookkeeping.
